\section{Naive set theory}\label{sec:naive_set_theory}

We will now describe some basics of \hyperref[def:naive_set_theory]{naive (axiomatic) set theory} without strictly adhering to the formalisms of \fullref{sec:formal_set_theory}.

\paragraph{Power sets}

\begin{definition}\label{def:subset}\mcite[10]{Мальцев1970АлгебраическиеСистемы}
  If every member of the \hyperref[def:set]{set} \( A \) is a member of \( B \), we say that \( A \) is a \term[bg=подмножество (\cite[144]{Тагамлицки1978ИнтегралноСмятане}), ru=подмножество, en=subset (\cite[def. 2.1.3]{Rosen2019DiscreteMathematics})]{subset} of \( B \) and write \( A \subseteq B \), and similarly we say that \( B \) is a \term[bg=подмножество (\cite[144]{Тагамлицки1978ИнтегралноСмятане}), ru=надмножество, en=superset (\cite[def. 2.1.3]{Rosen2019DiscreteMathematics})]{superset} of \( A \) and write \( B \supseteq A \).

  If \( A \subseteq B \) and \( A \neq B \), we say that the subset \( A \) is \term[ru=собственное (подмножество) (\cite[34]{БелоусовТкачёв2004ДискретнаяМатематика}), en=proper (subset) (\cite[126]{Rosen2019DiscreteMathematics})]{proper} or \term[ru=строгое (подмножество) (\cite[34]{БелоусовТкачёв2004ДискретнаяМатематика}), en=strict (subset) (\cite[403]{Müller2022HandbookOfDynamicsAndProbability})]{strict} of \( B \) and write \( A \subsetneq B \).

  \begin{thmenum}
    \thmitem{def:subset/formal} Formally, we can introduce the \hyperref[def:fol_signature/notation]{infix} binary \hyperref[def:fol_signature/predicate]{predicate symbol} \( {\synsubseteq} \) via the \hyperref[def:fol_definition/predicate]{first-order predicate definition}
    \begin{equation*}
      \qforall \syna \qforall \synb (\syna \synsubseteq \synb \syniff (\synx \synin \syna \synimplies \synx \synin \synb)).
    \end{equation*}
  \end{thmenum}
\end{definition}

\begin{remark}\label{rem:strict_subset_notation}
  Based on how we denote the \hyperref[def:partially_ordered_set/strict]{strict order} of \( {\leq} \) by \( {<} \), it is natural to use \( {\subset} \) for \hyperref[def:subset]{strict subsets}.

  Unfortunately, a lot of authors use \( A \subset B \) for the nonstrict relation, so we prefer the less ambiguous \( A \subsetneq B \).

  See \cref{tab:rem:strict_subset_notation} for a usage comparison across authors.

  \begin{table}
    \centering
    \begin{tabular}{l l l}
      \toprule
      Source                                                          & Nonstrict subset    & Strict subset        \\
      \midrule
      \cite[9]{Jech2006SetTheory}                                     & \( A \subset B \)   & ---                  \\
      \cite[def. 1.3]{Rudin1976AnalysisPrinciples}                    & \( A \subset B \)   & ---                  \\
      \cite[1]{Engelking1989GeneralTopology}                          & \( A \subset B \)   & ---                  \\
      \cite[2]{Kelley1975GeneralTopology}                             & \( A \subset B \)   & ---                  \\
      \cite[91]{Halmos1960NaiveSetTheory}                             & \( A \subset B \)   & ---                  \\
      \cite[1]{Willard2004GeneralTopology}                            & \( A \subset B \)   & ---                  \\
      \cite[7]{Зорич2019АнализЧасть1}                                 & \( A \subset B \)   & ---                  \\
      \cite[59]{ИльинСадовничийСендов1985АнализТом1}                  & \( A \subset B \)   & ---                  \\
      \cite[15]{КусраевКутателадзе2005БулевозначныйАнализ}            & \( A \subset B \)   & ---                  \\
      \cite[144]{Тагамлицки1978ИнтегралноСмятане}                     & \( A \subset B \)   & ---                  \\
      \cite[3]{Jacobson1985BasicAlgebraI}                             & \( A \subset B \)   & \( A \subsetneq B \) \\
      \cite[13]{ШеньВерещагин2020НачалаТеорииМножеств}                & \( A \subset B \)   & \( A \subsetneq B \) \\
      \cite[17]{Kunen2013SetTheory}                                   & \( A \subseteq B \) & ---                  \\
      \cite[159]{Enderton1977SetTheory}                               & \( A \subseteq B \) & ---                  \\
      \cite[3]{Harzheim2005OrderedSets}                               & \( A \subseteq B \) & \( A \subset B \)    \\
      \cite[7]{Savage1998ModelsOfComputation}                         & \( A \subseteq B \) & \( A \subset B \)    \\
      \cite[135]{Kleene2002Logic}                                     & \( A \subseteq B \) & \( A \subset B \)    \\
      \cite[503]{Müller2022HandbookOfDynamicsAndProbability}          & \( A \subseteq B \) & \( A \subset B \)    \\
      \cite[def. 2.1.3]{Rosen2019DiscreteMathematics}                 & \( A \subseteq B \) & \( A \subset B \)    \\
      \cite[def. 1.1]{Александров1977ОбщаяТопология}                  & \( A \subseteq B \) & \( A \subset B \)    \\
      \cite[def. 1.1]{АлександровКолмогоров1948ТеорияМножеств}        & \( A \subseteq B \) & \( A \subset B \)    \\
      \cite[10]{Мальцев1970АлгебраическиеСистемы}                     & \( A \subseteq B \) & \( A \subset B \)    \\
      \cite[33]{БелоусовТкачёв2004ДискретнаяМатематика}               & \( A \subseteq B \) & \( A \subset B \)    \\
      \cite[def. 3.1.14]{Tao2022AnalysisI}                            & \( A \subseteq B \) & \( A \subsetneq B \) \\
      \bottomrule
    \end{tabular}
    \caption{The table for \cref{rem:strict_subset_notation} discussing how different authors denote \hyperref[def:subset]{subsets}.}\label{tab:rem:strict_subset_notation}
  \end{table}
\end{remark}

\begin{definition}\label{def:power_set}\mcite[19]{Enderton1977SetTheory}
  For every set \( A \), we define the \term[ru=булеан (\cite[34]{БелоусовТкачёв2004ДискретнаяМатематика})]{power set} \( \pow(A) \) whose members are the \hyperref[def:subset]{subsets} of \( A \).

  \begin{thmenum}
    \thmitem{def:power_set/formal} Formally, in accordance with \cref{def:set_theory_definition}, we can use set-builder notation to introduce a unary \hyperref[def:fol_signature/function]{function symbol} \( {\syn\pow} \):
    \begin{equation*}
      \syn\pow(\syna) \coloneqq \set{ \synx \given \synx \synsubseteq \syna }.
    \end{equation*}

    We have used the subset predicate symbol \( {\synsubseteq} \) from \cref{def:subset}.
  \end{thmenum}
\end{definition}
\begin{comments}
  \item In \hyperref[def:zfc]{\logic{ZFC}}, the correctness of this definition relies on the \hyperref[def:zfc/power_set]{power set axiom} --- see \cref{thm:zfc_existence_theorems/power_set}.
\end{comments}

\paragraph{Sets of sets}

\begin{remark}\label{rem:sets_of_sets}
  In a \hyperref[def:axiomatic_set_theory/atoms]{pure set theory}, the members of sets are always sets.

  When speaking about numbers, functions, symbols or abstract structures, we regard their encoding as sets is a technicality. Thus, practically, speaking about sets of numbers or functions is not fundamentally different from speaking about informal collections of numbers or functions.

  On the other hand, we often have to deal with sets whose members we also regard as sets. To avoid repetition, it is convenient to call them \enquote{collections} or \enquote{families} of sets (\enquote{\textru{совокупность}} and \enquote{\textru{семейство}} in Russian). Both terms are introduced as a synonym for \enquote{set} in
  \cite[12]{Jech2006SetTheory},
  \cite[1]{Kelley1975GeneralTopology},
  \cite[1]{Willard2004GeneralTopology},
  \cite[27]{Rudin1976AnalysisPrinciples},
  \cite[2]{Folland1999RealAnalysis} and
  \cite[5]{Зорич2019АнализЧасть1}.

  We will avoid using the phrase \enquote{family of sets} because the word \enquote{family} can also refer to indexed family as defined in \cref{def:indexed_family}. This usage is introduced, for example,
  \cite[44]{Halmos1960NaiveSetTheory},
  \cite[81]{Bourbaki2004TheoryOfSets} and
  \cite[3]{Harzheim2005OrderedSets}.
\end{remark}

\begin{definition}\label{def:empty_set}\mcite[17]{Kunen2013SetTheory}
  The dual of the paradoxical \hyperref[def:universal_set]{universal set} is the \term{empty set} \( \varnothing \) with no members.

  \begin{thmenum}
    \thmitem{def:empty_set/formal} Formally, in accordance with \cref{def:set_theory_definition}, we can use set-builder notation to introduce a nullary \hyperref[def:fol_signature/function]{function symbol} \( {\syn\varnothing} \):
    \begin{equation*}
      \syn\varnothing \coloneqq \set{ \synx \given \synbot }.
    \end{equation*}
  \end{thmenum}
\end{definition}
\begin{comments}
  \item In \hyperref[def:zfc]{\logic{ZFC}}, the correctness of this definition relies on the \hyperref[rem:zfc_formulations/empty_set]{axiom of the empty set} or an equivalent --- see \cref{thm:zfc_existence_theorems/empty_set}.

  \item This is an \hyperref[def:universal_objects/initial]{initial object} in the \hyperref[def:category_of_small_sets]{category of small sets} --- see \cref{thm:universal_sets}.
\end{comments}

\begin{definition}\label{def:subsingleton_set}\mcite[112]{UnivalentFoundationsProgram2013HoTT}
  If all elements of a set are equal, we call it a \term{subsingleton set}.
\end{definition}
\begin{comments}
  \item The term \enquote{subsingleton} is suggested as a term for \hyperref[def:mere_proposition]{mere proposition} in \cite[112]{UnivalentFoundationsProgram2013HoTT}. We adapt its usage for axiomatic set theory.

  On the other hand, \enquote{singleton} is a more popular term, given as a definition of axiomatic set theory in
  \cite[10]{Halmos1960NaiveSetTheory},
  \cite[7]{Jech2006SetTheory},
  \cite[3]{Kelley1975GeneralTopology},
  \cite[131]{КановейЛюбецкий2010СовременнаяТеорияМножеств} and
  \cite[129]{КолмогоровДрагалин2006ЛогикаВведение}.

  \item An empty set has zero elements, a singleton has one unique element, while a subsingleton may be either the empty set or a singleton. We give this more abstract definition which avoids direct references to numbers.
\end{comments}

\begin{definition}\label{def:singleton_set}\mcite[10]{Halmos1960NaiveSetTheory}
  We call a \hyperref[def:empty_set]{nonempty} \hyperref[def:singleton_set]{subsingleton set} a \term[ru=синглет (\cite[131]{КановейЛюбецкий2010СовременнаяТеорияМножеств}); синглетон (\cite[129]{КолмогоровДрагалин2006ЛогикаВведение})]{singleton}.

  Via \hyperref[def:set_theory_definition/roster]{roster notation}, we can denote the singleton containing \( a \) by \( \set{ a } \).
\end{definition}
\begin{comments}
  \item Since a singleton set is a special case of an \hyperref[def:unordered_pair]{unordered pair}, in \hyperref[def:zfc]{\logic{ZFC}}, the correctness of this definition relies on the \hyperref[def:zfc/pairing]{axiom of pairing} --- see \cref{thm:zfc_existence_theorems/singleton}.

  \item The singleton set \( \set{ A } \), where \( A \) is a \hyperref[def:small_set]{small set}, is a \hyperref[def:universal_objects/terminal]{terminal object} in the \hyperref[def:category_of_small_sets]{category of small sets} --- see \cref{thm:universal_sets}.

  \item It is sometimes convenient, especially in connection with \hyperref[ch:geometry]{geometry} or \hyperref[def:function]{set-valued maps} (e.g. when dealing with \hyperref[def:net_limit_point]{limits of nets} or \hyperref[def:subdifferentials]{subdifferentials}), to not distinguish between singleton sets and their corresponding element.
\end{comments}

\begin{definition}\label{def:unordered_pair}\mcite[9]{Halmos1960NaiveSetTheory}
  Given sets \( a \) and \( b \), via \hyperref[def:set_theory_definition/roster]{roster notation} we can define the \term[ru=неупорядоченная пара (\cite[181]{Герасимов2014Вычислимость})]{unordered pair} \( \set{ a, b } \).

  We do not require \( a \) and \( b \) to be distinct. In fact, if they are equal, we obtain a \hyperref[def:singleton_set]{singleton set}.
\end{definition}
\begin{comments}
  \item In \hyperref[def:zfc]{\logic{ZFC}}, the correctness of this definition relies on the \hyperref[def:zfc/pairing]{axiom of pairing} --- see \cref{thm:zfc_existence_theorems/unordered_pair}.

  \item We only need unordered pairs as a technical tool for defining binary unions and intersections in this sections and Kuratowski pairs in \cref{def:kuratowski_pair}. They are not by themselves an important concept.
\end{comments}

\begin{proposition}\label{thm:unordered_pair_extensionality}
  For \hyperref[def:unordered_pair]{unordered pairs} of sets the following property holds:
  \begin{equation}\label{eq:thm:unordered_pair_extensionality}
    \set{ a, b } = \set{ c, d } \T{if and only if} (a = c \T{and} b = d) \T{or} (a = d \T{and} b = c).
  \end{equation}
\end{proposition}
\begin{comments}
  \item For \hyperref[def:kuratowski_pair]{Kuratowski pairs}, we have the stronger statement \cref{thm:unordered_pair_extensionality}. For indexed families, we have \cref{thm:indexed_family_extensionality}.
\end{comments}
\begin{proof}
  \SufficiencySubProof Suppose that \( \set{ a, b } = \set{ c, d } \).

  By the \hyperref[def:naive_set_theory/extensionality]{axiom of extensionality}, both sets have the same members. We will consider several cases:
  \begin{itemize}
    \item If \( a = b \), then both \( c \) and \( d \) must be equal to their common value, hence \( a = b = c = d \). In particular, we have both
    \begin{equation*}
      a = c \T{and} b = d
    \end{equation*}
    and
    \begin{equation*}
      a = d \T{and} b = c.
    \end{equation*}

    \item If \( a \neq b \) and \( a = c \), then \( b = c \) leads to a contradiction, so it remains for \( b = d \) to hold.

    \item Similarly, if \( a \neq b \) and \( a = d \), then \( b = c \).
  \end{itemize}

  \NecessitySubProof In the converse, we also do case analysis.
  \begin{itemize}
    \item If \( a = c \) and \( b = d \), then both \( a \) and \( b \) belong to \( \set{ c, d } \), hence \( \set{ a, b } \) is a subset of \( \set{ c, d } \). We obtain the converse inclusion analogously, thus \( \set{ a, b } = \set{ c, d } \).

    \item If \( a = d \) and \( b = c \), we proceed analogously.
  \end{itemize}
\end{proof}

\paragraph{Operations on sets}

\begin{definition}\label{def:basic_set_operations}
  We define the following operations:

  \begin{thmenum}
    \thmitem{def:basic_set_operations/union}\mcite[23]{Enderton1977SetTheory} We define the \term{union} of an arbitrary set \( \mscrA \) as
    \begin{equation*}
      \bigcup A \coloneqq \set{ x \given \qexists {A \in \mscrA} x \in A }.
    \end{equation*}

    The corresponding \hyperref[con:formula_defined_predicate]{formula-defined predicate} is
    \begin{equation*}\taglabel[\op{IsUnion}]{eq:def:basic_set_operations/union/predicate}
      \ref{eq:def:basic_set_operations/union/predicate}[\rho, \sigma] \coloneqq \qforall \synx \parens[\big]{ \synx \in \rho \syniff \qexists {\syny \in \sigma} \synx \in \syny }.
    \end{equation*}

    In particular, \( \bigcup \varnothing = \varnothing \).

    For two sets \( A \) and \( B \), we define their \term{binary union} as
    \begin{equation*}
      A \cup B \coloneqq \bigcup \set{ A, B } = \set{ x \given x \in A \T{or} x \in B }.
    \end{equation*}

    \thmitem{def:basic_set_operations/intersection}\mcite[24]{Enderton1977SetTheory} Dually to \hyperref[def:basic_set_operations/union]{unions}, we define the \term{intersection} of a nonempty set \( \mscrA \) as
    \begin{equation*}
      \bigcap \mscrA \coloneqq \set{ x \given \qforall {A \in \mscrA} x \in A }.
    \end{equation*}

    We also introduce the \hyperref[con:formula_defined_predicate]{formula-defined predicate}
    \begin{equation*}\taglabel[\op{IsIntersection}]{eq:def:basic_set_operations/intersection/predicate}
      \ref{eq:def:basic_set_operations/intersection/predicate}[\rho, \sigma] \coloneqq \qforall \synx \parens[\big]{ \synx \in \rho \syniff \qforall {\syny \in \sigma} \synx \in \syny }.
    \end{equation*}

    We leave \( \bigcap \varnothing \) undefined because it should be a \hyperref[def:extremal_points/top_and_bottom]{top element} in the \hyperref[thm:boolean_algebra_of_subsets]{Boolean algebra of all sets}, but the latter object is ambiguous and does not even exist in \logic{ZFC} --- see \cref{thm:zfc_existence_theorems/set_containing_all}.

    For two sets \( A \) and \( B \), we define their \term{binary intersection} as
    \begin{equation*}
      A \cap B \coloneqq \bigcap \set{ A, B } = \set{ x \given x \in A \T{and} x \in B }.
    \end{equation*}

    \thmitem{def:basic_set_operations/difference}\mcite[27]{Enderton1977SetTheory} We define the \term{difference} of \( A \) and \( B \) as
    \begin{equation*}
      A \setminus B \coloneqq \set{ x \in A \given x \notin B }.
    \end{equation*}

    We define the \hyperref[con:formula_defined_predicate]{formula-defined predicate}
    \begin{equation*}\taglabel[\op{IsDifference}]{eq:def:basic_set_operations/difference/predicate}
      \ref{eq:def:basic_set_operations/difference/predicate}[\rho, \sigma, \tau] \coloneqq \qforall \synx \parens[\big]{ \synx \in \rho \syniff (\synx \in \sigma \synwedge \synneg(\synx \in \tau)) }.
    \end{equation*}

    \thmitem{def:basic_set_operations/symmetric_difference}\mcite[9]{Jech2006SetTheory} We define the \term{symmetric difference} of two sets \( A \) and \( B \) as
    \begin{equation}
      A \symdiff B \coloneqq \set{ x \in A \given x \notin B }.
    \end{equation}

    We also define the \hyperref[con:formula_defined_predicate]{formula-defined predicate}
    \begin{equation*}\taglabel[\op{IsSymDiff}]{eq:def:basic_set_operations/symmetric_difference/predicate}
      \ref{eq:def:basic_set_operations/symmetric_difference/predicate}[\rho, \sigma, \tau] \coloneqq \ref{eq:def:basic_set_operations/difference/predicate}[\rho, \sigma, \tau] \synvee \ref{eq:def:basic_set_operations/difference/predicate}[\rho, \tau, \sigma].
    \end{equation*}
  \end{thmenum}
\end{definition}

\begin{proposition}\label{thm:union_intersection_distributivity}
  \hyperref[def:basic_set_operations/union]{Set unions} and \hyperref[def:basic_set_operations/intersection]{set intersections} distribute over themselves and over each other. That is,
  \begin{align}
    X \cup \bigcup \mscrA &= \bigcup \set{ A \cup X \given A \in \mscrA }, \label{eq:thm:union_intersection_distributivity/union_over_union} \\
    X \cap \bigcup \mscrA &= \bigcup \set{ A \cap X \given A \in \mscrA }, \label{eq:thm:union_intersection_distributivity/intersection_over_union} \\
    X \cup \bigcap \mscrA &= \bigcap \set{ A \cup X \given A \in \mscrA }, \label{eq:thm:union_intersection_distributivity/union_over_intersection} \\
    X \cap \bigcap \mscrA &= \bigcap \set{ A \cap X \given A \in \mscrA }. \label{eq:thm:union_intersection_distributivity/intersection_over_intersection}
  \end{align}
\end{proposition}
\begin{proof}
  \SubProofOf{eq:thm:union_intersection_distributivity/union_over_union} Trivial.

  \SubProofOf{eq:thm:union_intersection_distributivity/union_over_intersection}
  \begin{align*}
    X \cap \bigcup \mscrA
    &=
    \set{ x \given x \in X \T{and} \qexists{A \in \mscrA} x \in A }
    = \\ &=
    \set{ x \given \qexists{A \in \mscrA} (x \in A \T{and} x \in X) }
    = \\ &=
    \bigcup \set{ A \cup X \given A \in \mscrA }.
  \end{align*}

  \SubProofOf{eq:thm:union_intersection_distributivity/intersection_over_union} Analogous.

  \SubProofOf{eq:thm:union_intersection_distributivity/intersection_over_intersection} Trivial.
\end{proof}

\begin{proposition}\label{thm:set_difference}
  \hyperref[def:basic_set_operations/difference]{Set difference} has the following basic properties:
  \begin{thmenum}
    \thmitem{thm:set_difference/intersection} If \( A \) and \( B \) are subsets of \( C \), then \( A \setminus B = A \cap (C \setminus B) \).

    \thmitem{thm:set_difference/superset} For any two sets \( A \) and \( B \), we have \( A \setminus B = A \setminus (A \cap B) \).

    \thmitem{thm:set_difference/double_difference} If \( A \subseteq B \), then \( B \setminus (B \setminus A) = A \)

    \thmitem{thm:set_difference/complement} If \( A \subseteq B \subseteq C \), then \( C \setminus B \subseteq C \setminus A \).

    If \( A \neq B \), then \( C \setminus B \neq C \setminus A \).
  \end{thmenum}
\end{proposition}
\begin{proof}
  \SubProofOf{thm:set_difference/intersection} Since \( a \in A \) implies \( a \in C \), we have
  \begin{align*}
    A \setminus B
    &=
    \set{ x \in A \given x \notin B }
    = \\ &=
    \set{ x \in A \given x \in C \T{and} x \notin B }
    = \\ &=
    A \cap (C \setminus B).
  \end{align*}

  \SubProofOf{thm:set_difference/superset} Follows from \cref{thm:set_difference/intersection}.

  \SubProofOf{thm:set_difference/double_difference} By \hyperref[thm:propositional_semantic_dne]{double negation elimination},
  \begin{align*}
    B \setminus (B \setminus A)
    &=
    \set[\big]{ x \in B \given x \notin \set{ x \in B \given x \notin A } }
    = \\ &=
    \set{ x \in B \given x \in A }
    = \\ &=
    A.
  \end{align*}

  \SubProofOf{thm:set_difference/complement} Any element in \( A \) is also in \( B \), thus, by contraposition, every element not in \( B \) is not in \( A \).
\end{proof}

\begin{proposition}\label{thm:boolean_algebra_of_subsets}
  Let \( X \) be an arbitrary set. Then the \hyperref[def:power_set]{power set} \( \pow(X) \) endowed with the \hyperref[def:subset]{inclusion} partial order \( \subseteq \) is a \hyperref[def:complete_lattice]{complete} \hyperref[def:boolean_algebra]{Boolean algebra}\fnote{We will not make use here of the full apparatus of Boolean algebras. Instead, we regard this statement as a shorthand that would otherwise require enumerating a lot of properties.}.

  Explicitly:
  \begin{thmenum}
    \thmitem{thm:boolean_algebra_of_subsets/join} The \hyperref[def:lattice]{join} of an arbitrary family \( \mscrA \) of subsets of \( X \) is simply the \hyperref[def:basic_set_operations/union]{union} \( \bigcap \mscrA \).

    \thmitem{thm:boolean_algebra_of_subsets/top} The \hyperref[def:extremal_points/top_and_bottom]{top element} is the set \( X \) itself.

    \thmitem{thm:boolean_algebra_of_subsets/meet} The \hyperref[def:lattice]{meet} of an arbitrary family \( \mscrA \) of sets is simply the \hyperref[def:basic_set_operations/intersection]{intersection} \( \bigcup \mscrA \). Unlike for a general family of sets, we have no problem defining the intersection of an empty set to be the top element \( X \).

    \thmitem{thm:boolean_algebra_of_subsets/bottom} The \hyperref[def:extremal_points/top_and_bottom]{bottom element} is the empty set.

    \thmitem{thm:boolean_algebra_of_subsets/complement} The \hyperref[def:boolean_algebra]{complement} \( A^\complement \) of the subset \( A \) is the \hyperref[def:basic_set_operations/difference]{difference} \( X \setminus A \).
  \end{thmenum}

  \begin{figure}
    \hfill
    \includegraphics[page=1]{output/thm__boolean_algebra_of_subsets}
    \hfill\hfill
    \caption{The \hyperref[def:hasse_diagram]{Hasse diagram} of \( \pow(\set{ A, B }) \) with respect to \hyperref[def:subset]{set inclusion}}
    \label{fig:thm:boolean_algebra_of_subsets}
  \end{figure}
\end{proposition}
\begin{proof}
  \SubProofOf{thm:boolean_algebra_of_subsets/join} The union of \( \mscrA \) exists by \cref{thm:zfc_existence_theorems/arbitrary_union}, and it is itself a subset of \( \mscrA \). Every set in \( \mscrA \) is contained in \( \bigcup \mscrA \), hence it is indeed a join.

  \SubProofOf{thm:boolean_algebra_of_subsets/top} Clearly \( X \) contains every subset of \( X \).

  \SubProofOf{thm:boolean_algebra_of_subsets/join} The intersection of \( \mscrA \) exists by \cref{thm:zfc_existence_theorems/arbitrary_union}, and it is itself a subset of \( \mscrA \). Every set in \( \mscrA \) contains \( \bigcap \mscrA \), hence it is indeed a meet.

  \SubProofOf{thm:boolean_algebra_of_subsets/bottom} The empty set is contained in every set, in particular in every subset of \( A \).

  \SubProofOf[def:distributive_lattice]{distributivity} Follows from \cref{thm:union_intersection_distributivity}.

  \SubProofOf{thm:boolean_algebra_of_subsets/complement} The operation \( A^\complement \) is well-defined for each subset \( A \) of \( X \) due to \cref{thm:zfc_existence_theorems/difference}.

  By definition
  \begin{equation*}
    A \synvee A^\complement
    =
    A \cup (X \setminus A)
    =
    X
  \end{equation*}
  and
  \begin{equation*}
    A \synvee A^\complement
    =
    A \cup (X \setminus A)
    =
    X,
  \end{equation*}
  hence \( A^\complement \) is indeed the complement of \( A \).

  Therefore, \( \pow(X) \) is a Boolean algebra.
\end{proof}

\begin{theorem}[De Morgan's laws for sets]\label{thm:de_morgans_laws_for_sets}
  The following hold for any sets:
  \begin{align}
    X \setminus \bigcup \mscrA = \bigcap \set{ X \setminus A \given A \in \mscrA } \label{eq:thm:de_morgans_laws_for_sets/complement_of_union} \\
    X \setminus \bigcap \mscrA = \bigcup \set{ X \setminus A \given A \in \mscrA } \label{eq:thm:de_morgans_laws_for_sets/complement_of_intersection}
  \end{align}
\end{theorem}
\begin{comments}
  \item This is one of several variants of De Morgan's laws presented here --- see \cref{rem:de_morgans_laws/sets}.
\end{comments}
\begin{proof}
  \SubProof{Direct proof} We have
  \begin{align*}
    X \setminus \bigcup \mscrA
    &=
    \set{ x \given x \in X \T{and} \synneg \qexists{A \in \mscrA} x \in A }
    \reloset {\eqref{eq:thm:fol_quantifier_duality}} = \\ &=
    \set{ x \given x \in X \T{and} \qforall{A \in \mscrA} {x \notin A} }
    = \\ &=
    \set{ x \given \qforall{A \in \mscrA} (x \in X \T{and} x \notin A) }
    = \\ &=
    \bigcap \set{ X \setminus A \given A \in \mscrA }
  \end{align*}
  and
  \begin{align*}
    X \setminus \bigcap \mscrA
    &=
    \set{ x \given x \in X \T{and} \synneg \qforall{A \in \mscrA} x \in A }
    \reloset {\eqref{eq:thm:fol_quantifier_duality}} = \\ &=
    \set{ x \given x \in X \T{and} \qexists{A \in \mscrA} {x \notin A} }
    = \\ &=
    \set{ x \given \qforall{A \in \mscrA} (x \in X \T{and} x \notin A) }
    = \\ &=
    \bigcup \set{ X \setminus A \given A \in \mscrA }.
  \end{align*}

  \SubProof{Consequence of \fullref{thm:de_morgans_laws}}
  All subsets of \( X \) form a Boolean algebra as a consequence of \cref{thm:boolean_algebra_of_subsets}. Then
  \begin{equation*}
    \bigcup_{k \in \mscrK} (X \setminus A_k)
    \reloset {\ref{thm:set_difference/superset}} =
    \bigcup_{k \in \mscrK} (X \cap A_k)^\complement
    \reloset {\eqref{eq:thm:de_morgans_laws/complement_of_join}} =
    \parens*{ \bigcap_{k \in \mscrK} (X \cap A_k) }^\complement
    =
    \parens*{ X \cap \bigcap_{k \in \mscrK} A_k }^\complement
    \reloset {\ref{thm:set_difference/superset}} =
    X \setminus \parens*{ \bigcap_{k \in \mscrK} A_k }.
  \end{equation*}

  The other identity requires distributivity:
  \begin{equation*}
    \bigcap_{k \in \mscrK} (X \setminus A_k)
    \reloset {\ref{thm:set_difference/superset}} =
    \bigcap_{k \in \mscrK} (X \cap A_k)^\complement
    \reloset {\eqref{eq:thm:de_morgans_laws/complement_of_join}} =
    \parens*{ \bigcup_{k \in \mscrK} (X \cap A_k) }^\complement
    \reloset {\eqref{eq:thm:union_intersection_distributivity/intersection_over_union}} =
    \parens*{ X \cap \bigcup_{k \in \mscrK} A_k }^\complement
    \reloset {\ref{thm:set_difference/superset}} =
    X \setminus \parens*{ \bigcup_{k \in \mscrK} A_k }.
  \end{equation*}
\end{proof}

\begin{remark}\label{rem:inductive_sets}
  We are now interested in the inductive axiom schema \eqref{eq:def:peano_arithmetic/theory/PA3}, which will allow us to define, among other things, natural numbers. Induction is an important proof technique that is discussed in detail in our proof of \cref{thm:peano_arithmetic_predecessor_existence}. There are more general forms of induction than \eqref{eq:def:peano_arithmetic/theory/PA3} like \fullref{thm:well_founded_induction} and \fullref{thm:well_founded_induction}. They do, however, require concepts which in turn depend on the existence of natural numbers within set theory. As a consequence, we cannot prove \eqref{eq:def:peano_arithmetic/theory/PA3} via \fullref{thm:well_founded_induction}.

  We will introduce the concept of inductive sets in \cref{def:inductive_set} and prove in \fullref{thm:omega_induction} that a special inductive set \hyperref[thm:smallest_inductive_set_existence]{\( \omega \)}, which will be the universe of our model of \( \BbbN \), allows performing inductive proofs. The technique that allows us to perform inductive proofs on \( \omega \) can be seen in our proof of \cref{thm:omega_is_transitive}. \Fullref{thm:omega_induction} will allow us to define natural numbers without relying on metalogical induction along the way. See our proof of \fullref{thm:omega_induction} for a description of natural number induction within set theory and \cref{rem:def:peano_arithmetic_term_model} for a further discussion of the use of natural numbers in the metatheory and in the object logic.

  We also introduce \term{recursion} in parallel as a technique for constructing objects. See \fullref{thm:omega_recursion}.
\end{remark}

\begin{definition}\label{def:ordinal_successor}\mcite[68]{Enderton1977SetTheory}
  We define the \term[ru=операция следования (\cite[1.1.11]{КолмогоровДрагалин2006ЛогикаВведение})]{successor} \( \op{sc}(A) \) of a set \( A \) as the set
  \begin{equation*}
    \op{sc}(A) \coloneqq A \cup \set{ A }.
  \end{equation*}

  The following \hyperref[con:formula_defined_predicate]{formula-defined predicate}
  \begin{equation*}\taglabel[\op{IsSucc}]{eq:def:ordinal_successor/predicate}
    \ref{eq:def:ordinal_successor/predicate}[\rho, \sigma] \coloneqq \qforall \synx \parens[\big]{ \synx \in \rho \syniff (\synx \in \sigma \synvee \synx = \sigma) },
  \end{equation*}
  which states that \( \rho \) is the successor of \( \sigma \), will be useful for us when working with \hyperref[def:inductive_set]{inductive sets}.
\end{definition}
\begin{comments}
  \item We also call it the \term{ordinal successor} operation since it is an important concept in the \hyperref[sec:ordinals]{theory of ordinals}. See \cref{rem:def:ordinal_successor} for an example of how it naturally arises. It should be distinguished from successor cardinals defined in \cref{def:successor_and_limit_cardinal}.
\end{comments}

\begin{definition}\label{def:inductive_set}\mcite[68]{Enderton1977SetTheory}
  A set is called \term{inductive} if contains the empty set and is closed under the \hyperref[def:ordinal_successor]{successor operator}.

  We introduce the following \hyperref[con:formula_defined_predicate]{formula-defined predicate}
  \begin{equation*}\taglabel[\op{IsInductive}]{eq:def:inductive_set/predicate}
    \ref{eq:def:inductive_set/predicate}[\sigma] \coloneqq
      \parens[\big]{ \qexists {\synx \in \sigma} \ref{eq:def:empty_set/predicate}[\synx] }
      \synwedge
      \parens[\big]{ \qforall {\synx \in \sigma} \qexists {\syny \in \sigma} \ref{eq:def:ordinal_successor/predicate}[\syny, \synx] }.
  \end{equation*}
\end{definition}

\begin{proposition}\label{thm:smallest_inductive_set_existence}
  There is a smallest (with respect to set inclusion) \hyperref[def:inductive_set]{inductive set}, which we denote by \( \omega \).
\end{proposition}
\begin{proof}
  We cannot directly define \( \omega \) as the intersection of all inductive sets since, in \logic{ZFC}, we want to avoid unrestricted comprehension. Fortunately, the existence of at least one inductive set is justified by the \hyperref[def:zfc/infinity]{axiom of infinity} (or by taking the entire universe in naive set theory).

  Let \( A \) be an inductive set. We can now define \( \omega \) via restricted comprehension:
  \begin{equation*}
    \omega \coloneqq \set{ x \in A \given x \T{belongs to every inductive set} }.
  \end{equation*}

  To see that \( \omega \) is itself inductive, note that \( \varnothing \in \omega \) and that if \( x \in \omega \), then it also belongs to all inductive sets and hence \( \op{sc}(x) \) also belongs to all inductive sets, proving \( \op{sc}(x) \in \omega \).
\end{proof}

\begin{theorem}[Induction via inductive sets]\label{thm:omega_induction}\mcite[thm. 6.2.17]{Hinman2005Logic}
  We can perform induction on the \hyperref[thm:smallest_inductive_set_existence]{smallest inductive set} \( \omega \). We can prove that \( \omega \) is a subset of some set \( A \) by proving the following:
  \begin{itemize}
    \item \( A \) contains \( \varnothing \).
    \item If \( A \) contains \( n \), it also contains \( \op{sc}(n) \).
  \end{itemize}

  More precisely, the following is a theorem of both naive set theory and \hyperref[def:zfc]{\logic{ZF}} (\( \omega \) satisfies \( \sigma \) and \( \synx \) plays the role of \( A \) above):
  \footnotesize
  \begin{equation}\label{eq:thm:omega_induction}
    \qexists \sigma
    \qforall \synx
    \parens[\Bigg]
      {
        \parens[\Bigg]
          {
            \parens[\big]
            {
              \underbrace{ \qexists {\syny \in \synx} \ref{eq:def:empty_set/predicate}[\syny] }_{\mathclap{\T{base case}}}
            }
            \synwedge
            \qforall \syny \parens[\big]
              {
                \overbrace
                  {
                    \underbrace{ \syny \in \synx }_{\mathclap{\substack{\T{inductive} \\ \T{hypothesis}}}}
                    \synimplies
                    \underbrace
                      {
                        \qexists {\synz \in \synx} \ref{eq:def:ordinal_successor/predicate}[\synz, \syny]
                      }_{\mathclap{\substack{\T{inductive step} \\ \T{conclusion}}}}
                  }^{\T{inductive step}}
              }
          }
        \synimplies
        \underbrace{ \ref{eq:def:subset/predicate}[\sigma, \synx] }_{\T{conclusion}}
      }
  \end{equation}
  \normalsize
\end{theorem}
\begin{comments}
  \item This is an analog of \eqref{eq:def:peano_arithmetic/theory/PA3} and is actually used in \cref{thm:omega_is_model_of_pa} to prove that \( \omega \) is a model of \hyperref[def:peano_arithmetic]{\logic{PA}}. Instead of an entire theorem schema, however, for this theorem it is sufficient to use one single formula. The more general induction principles that use theorem schemas cannot be proved without natural numbers, which are a model of \logic{PA} by virtue of this theorem.
\end{comments}
\begin{proof}
  The antecedent of (the inner formula in) \eqref{eq:thm:omega_induction} is a restatement of the formula-defined predicate \( \ref{eq:def:inductive_set/predicate}[\synx] \). The situation resembles the \hyperref[eq:def:zfc/infinity]{axiom of infinity}, but, instead of existence of an inductive set \( \synx \), it states the existence of a set \( \sigma \) such that if \( \synx \) is an inductive set, then \( \sigma \) is a subset of \( \synx \) (if we restrict \( \syny \) to range only over members of \( \sigma \), then we would obtain equality of \( \synx \) and \( \sigma \) instead). In other words, we have reduced the verification of \eqref{eq:thm:omega_induction} to showing that there exists a minimal inductive set in both naive set theory and \logic{ZF}.

  We have already proved in \cref{thm:smallest_inductive_set_existence} that our fixed model \( \mscrV = (V, I) \) of set theory has a minimal inductive set \( \omega \). Thus, for any \hyperref[def:fol_variable_assignment]{variable assignment} \( v \), the modified assignment \( v_{\sigma \mapsto \omega} \) satisfies \eqref{eq:thm:omega_induction} with the outer existential quantifier removed. Hence, \( v \) satisfies the entire formula \eqref{eq:thm:omega_induction}.

  Both the assignment \( v \) and the model \( \mscrV \) were arbitrary, therefore we can conclude that \eqref{eq:thm:omega_induction} is a theorem of both naive set theory and \logic{ZF}.
\end{proof}

\begin{definition}\label{def:transitive_set}\mcite[def. 6.1.3(i)]{Hinman2005Logic}
  We say that a set \( A \) is \term{transitive} if from \( B \in A \) it follows that \( B \subseteq A \).

  We introduce the following \hyperref[con:formula_defined_predicate]{formula-defined predicate}:
  \begin{equation*}\taglabel[\op{IsSetTransitive}]{eq:def:transitive_set/predicate}
    \ref{eq:def:transitive_set/predicate}[\sigma] \coloneqq \qforall {\synx \in \sigma} \qforall {\syny \in \synx} {\syny \in \sigma}
  \end{equation*}
\end{definition}
\begin{comments}
  \item See \cref{rem:ordinal_definition} for a discussion of the motivation and terminology of transitive sets and \cref{def:set_theory_model/transitive} for their importance.
\end{comments}

\begin{proposition}\label{thm:omega_is_transitive}
  The set \( \omega \) is transitive and every member of \( \omega \) is transitive.
\end{proposition}
\begin{comments}
  \item This proof demonstrates how \fullref{thm:omega_induction} is used. From now on, we will not be as explicit about the use of induction on \( \omega \).
\end{comments}
\begin{proof}
  \SubProof{Proof that all members of \( \omega \) are transitive} In order to demonstrate how \fullref{thm:omega_induction} works in practice, we will use inductive sets directly. Let \( T \subseteq \omega \) be the subset of all transitive members of \( \omega \). We will show that \( T \) is inductive.

  Clearly \( \varnothing \in T \)  because every member of \( \varnothing \) vacuously is a subset of \( \varnothing \).

  Now suppose that \( n \in T \) and let \( m \in \op{sc}(n) = n \cup \set{ n } \). If \( m = n \), then \( m \in \op{sc}(n) \) by definition of the successor operation. If \( m \in n \), then \( m \subseteq n \) by the inductive hypothesis and hence also \( m \subseteq n \cup \set{ n } = \op{sc}(n) \). Thus, \( \op{sc}(n) \) is also transitive.

  We have shown that \( T \) is inductive. Therefore, \( \omega \subseteq T \), i.e. every member of \( \omega \) is transitive.

  \SubProof{Proof that \( \omega \) is transitive} We will show that, for all members \( n \) of \( \omega \), we have \( n \subseteq \omega \).

  The case \( n = \varnothing \) is again trivial.

  Now suppose that \( n \subseteq \omega \) and let \( m \in \op{sc}(n) \). If \( m = n \), clearly \( m \subseteq \omega \). If \( m \in n \), then \( m \subseteq n \) and, since \( n \subseteq \omega \), we have \( m \subseteq \omega \) by transitivity of \( \subseteq \).

  Therefore, \( \omega \) is transitive.
\end{proof}

\begin{lemma}\label{thm:members_of_omega_do_not_contain_themselves}
  No element of \( \omega \) is a member of itself.
\end{lemma}
\begin{proof}
  We will again use \fullref{thm:omega_induction}.

  By definition, \( \omega \notin \omega \). Now suppose that \( n \notin n \) for some \( n \in \omega \).

  Aiming at a contradiction, suppose that \( \op{sc}(n) \in \op{sc}(n) \). The assumption that \( n = \op{sc}(n) = n \cup \set{ n } \) implies that \( n = \set{ n } \). The assumption that \( \op{sc}(n) \in n \) implies that \( n \in n \) since \( \op{sc}(n) \) is transitive by \cref{thm:omega_is_transitive}. In both cases we have \( n \in n \), which contradicts our inductive hypothesis. This contradiction shows that \( \op{sc}(n) \notin \op{sc}(n) \).

  \Fullref{thm:omega_induction} allows us to conclude that no member of \( \omega \) contains itself.
\end{proof}

\begin{theorem}\label{thm:omega_is_model_of_pa_without_operations}
  The \hyperref[thm:smallest_inductive_set_existence]{smallest inductive set} \( \omega \) satisfies the axioms \eqref{eq:def:peano_arithmetic/theory/PA1}-\eqref{eq:def:peano_arithmetic/theory/PA3} from \hyperref[def:peano_arithmetic]{Peano arithmetic} with the following interpretation:
  \begin{thmenum}
    \thmitem{thm:omega_is_model_of_pa_without_operations/zero} \hyperref[def:peano_arithmetic/signature/zero]{Zero} is interpreted as \( \varnothing \).

    \thmitem{thm:omega_is_model_of_pa_without_operations/succ} The \hyperref[def:peano_arithmetic/signature/succ]{successor} operation \( s \) is interpreted as \( \op{sc} \).
  \end{thmenum}

  We will generalize this theorem to \cref{thm:omega_is_model_of_pa} after we are able to define the arithmetic operations in \( \omega \).
\end{theorem}
\begin{proof}
  \SubProofOf{eq:def:peano_arithmetic/theory/PA1} Let \( n, m \in \omega \) and suppose that \( \op{sc}(n) = \op{sc}(m) \). If \( n = m \), there is nothing to prove.

  Suppose that \( n \neq m \). Thus, since
  \begin{equation*}
    n \cup \set{ n } = \op{sc}(n) = \op{sc}(m) = m \cup \set{ m },
  \end{equation*}
  we have both \( n \in m \) and \( m \in n \).

  \Cref{thm:omega_is_transitive} implies that \( n \) is transitive and hence \( n \in n \), which contradicts \cref{thm:members_of_omega_do_not_contain_themselves}.

  The obtained contradiction shows that \( n = m \).

  \SubProofOf{eq:def:peano_arithmetic/theory/PA2} Suppose that \( \varnothing \) has a predecessor \( n \in \omega \). Then
  \begin{equation*}
    \varnothing = \op{sc}(n) = n \cup \set{ n },
  \end{equation*}
  which implies that \( n \in \varnothing \). But this contradicts the \hyperref[def:empty_set]{definition of \( \varnothing \)}.

  Therefore, \( \varnothing \) has no predecessor.

  \SubProofOf{eq:def:peano_arithmetic/theory/PA3} It follows from \fullref{thm:omega_induction} that \eqref{eq:thm:omega_induction} is a theorem of \logic{ZF}. Let \( \mscrV = (V, I) \) be our ambient \hyperref[def:set_theory_model/standard]{standard} \hyperref[def:set_theory_model/transitive]{transitive} model of \logic{ZFC}.

  Fix any \hyperref[def:fol_variable_assignment]{variable assignment} \( v \). As in our proof of \fullref{thm:omega_induction}, we consider the modified assignment \( v_{\tau \mapsto \omega} \) that \enquote{eliminates} the outer existential quantifier in \eqref{eq:thm:omega_induction}.

  To show that \eqref{eq:thm:omega_induction} really corresponds to \eqref{eq:def:peano_arithmetic/theory/PA3} (and hence that \( \omega \) satisfies \eqref{eq:def:peano_arithmetic/theory/PA3}), fix some formula \( \varphi \) of \hyperref[def:peano_arithmetic]{Peano arithmetic} (\emph{not \logic{ZFC}!}) and suppose that \( \synx, \synz_1, \ldots, \synz_n \) are all of its free variables. Fix also some parameter values \( u_1, \ldots, u_n \in \omega \) and, as in \cref{def:set_theory_definition}, define the set
  \begin{equation*}
    A \coloneqq \set{ x \in \omega \given \Bracks{\varphi}(x, u_1, \ldots, x_n) }.
  \end{equation*}

  Since \eqref{eq:thm:omega_induction} is satisfied by \( v \), the inner formula in \eqref{eq:thm:omega_induction} (without the quantifiers over \( \tau \) and \( \sigma \)) is satisfied by \( v_{\tau \mapsto \omega, \sigma \mapsto A} \).

  Since our choice of parameters \( u_1, \ldots, u_n \) was arbitrary, we can conclude that the universal closure of \eqref{eq:def:peano_arithmetic/theory/PA3} is satisfied by \( \omega \) for every formula \( \varphi \) of \logic{PA}.
\end{proof}

\begin{remark}\label{rem:set_theory_natural_numbers_without_operations}
  Due to \cref{thm:omega_is_model_of_pa_without_operations}, we will henceforth identify the \hyperref[thm:smallest_inductive_set_existence]{smallest inductive set} \( \omega \) with the set \( \BbbN \) of \hyperref[def:natural_numbers]{natural numbers}.

  We are not yet able to add or multiply natural numbers, nor rely on their well-foundedness, however for all other purposes we are able to utilize them.

  Since the ordering in \cref{def:peano_arithmetic_ordering} is defined via addition, we must define some other ordering. Luckily, as we shall see in \fullref{sec:ordinals}, \( n < m \) corresponds to \( n \in m \). In particular, the members of \( m \) are ordered.

  As a consequence, every natural number \( n \) equals the set of all smaller natural numbers by the \hyperref[def:naive_set_theory/extensionality]{axiom of extensionality}. It is conventional to write \( \set{ 0, 1, \ldots, n } \) rather than \( \set{ m \given m \in \op{sc}(n) } \). The former notation will be fully justified in \fullref{sec:ordinals}.

  Until then, natural numbers will be useful, for example, in \cref{def:sequence}.
\end{remark}
